% Incorporación ampliativa: CG-006, SR-032--SR-035 del inventario REV05.
% Residencia: inmediatamente después de la construcción de R36, antes de G9.
% Propietario computacional: 16_CIERRE_GLOBAL_HMT_MD_2026-07-22/
% 01_SELECTOR_CONSTANTES/selector_interno_constantes.py, líneas 31--63,
% 131--180, 185--240 y 267--321.
% Recuperación y cambio de representación: output/REVISION_EDITORIAL_HMT_MD_
% 20260905/03_PRUEBAS/verificar_generacion_infinita_nonadica.py,
% líneas 54--77, 651--723 y 974--1147.
% Exposición de origen: ampliaciones_sucesoras_20260824/parte_i_ii/owners/
% U008_orbitas_elevacion_r36.tex, líneas 210--358.
% Estatuto: RESULTADO_RECUPERADO; exposición de enumeración y naturalidad
% reunida sin promover la concordancia en R36 a equivalencia global.

\section[Selection and transport of the first boundary]{Combinatorial selection and representation transport of the first joint boundary}
\label{r36c:sec:frontera}

The first joint boundary receives the blocks already produced by the APP seeds, tritic reduction and TPK lifting calendar. Its construction preserves the distinction between three operations: determining admissible distributions, orienting their two lateral channels and transporting incidence to the exceptional coordinate system. The present account develops these operations on the same finite state. Names of real coordinates designate their subsequent readers; the following selection conditions act on ternary words, transport matrices and integer incidences.

\subsection{Regional state preceding the boundary}

Let $\mathcal W\subset\mathbb F_3^6$ be the repertoire of $243$ visible words obtained from the APP--TRIT--TPK emission census. For a row vector $s\in\mathcal W$, define
\[
 b_0(s)=s,\quad b_1(s)=sL_0,\quad b_2(s)=sL_0^2,\quad
 b_3(s)=sL_0^2L_1,\quad b_4(s)=sL_0^2L_1^2,
\]
with operations in $\mathbb F_3$. The matrices previously derived from the regional registers are
\[
L_0=\begin{pmatrix}
2&2&2&1&2&1\\2&1&2&2&1&1\\1&0&0&1&0&2\\
0&0&1&1&1&0\\2&2&2&0&2&1\\2&1&2&1&0&0
\end{pmatrix},\qquad
L_1=\begin{pmatrix}
2&1&2&0&2&2\\1&2&1&1&2&0\\1&2&1&1&1&2\\
0&1&0&0&2&1\\2&0&2&1&0&1\\0&2&2&2&1&1
\end{pmatrix}.
\]
The preceding regional selector yields
\[
s_{\rm cl}=010211,\qquad s_{\rm pr}=201101,
\]
through the conditions $b_2(s_{\rm cl})=b_0(s_{\rm cl})$, $b_4(s_{\rm pr})=b_2(s_{\rm pr})$ and $b_1(s_{\rm cl})=b_3(s_{\rm pr})$. At this stage, the current blocks are
\[
b_4(s_{\rm cl})=110221,\qquad
b_4(s_{\rm pr})=102011.
\]
Concatenation preserves the five blocks of each register; vector $b_4$ is the six-coordinate input received by the boundary relation.

To express the recovered relation faithfully, the incidence matrix is used in its original representation,
\begin{equation}
A^{\rm o}=\begin{pmatrix}
0&1&1&1&1&1\\1&0&1&1&2&2\\1&1&0&2&1&2\\
1&1&2&0&2&1\\1&2&1&2&0&1\\1&2&2&1&1&0
\end{pmatrix}.
\label{r36c:eq:ao}
\end{equation}
We verify $(A^{\rm o})^2=-I$ in $M_6(\mathbb F_3)$; therefore, $(A^{\rm o})^{-1}=-A^{\rm o}$. This representation, together with its vectors, will later be transported to the one used by the exceptional reader.

\subsection{Compatibility equation and column saturation}

The boundary axis is fixed by $z=b_4(s_{\rm pr})=102011$. For each $s\in\mathcal W$, define
\begin{align}
m(s)&=b_4(s)A^{\rm o}L_1^3,\\
q(s)&=m(s)+z,\\
\widetilde q(s)&=s(A^{\rm o})^{-1}L_0^3(A^{\rm o})^3.
\label{r36c:eq:compat-datos}
\end{align}
The compatibility condition is $q(s)=\widetilde q(s)$. The lateral rows $x,y\in\mathbb F_3^6$ satisfy $x+y=m(s)$; the ordered boundary is $B=(x,y,z)^{\mathsf T}$.

To form remainder and carry, the integer representatives $0,1,2$ of each trit are used. If $S_j=x_j+y_j+z_j$, require
\[
q_j=S_j\bmod3,\qquad c_j=\frac{S_j-q_j}{3}=1,
\]
and column occupation is given by
\begin{equation}
\mathbf1_{\{x_j\ne0\}}+\mathbf1_{\{y_j\ne0\}}+\mathbf1_{\{z_j\ne0\}}
=2+\mathbf1_{\{z_j\ne0,\ m_j(s)\ne2\}}.
\label{r36c:eq:ocupacion}
\end{equation}
Dual neutrality completes the condition through
\begin{equation}
BA^{\rm o}\mathbf1=0\quad\hbox{in }\mathbb F_3^3.
\label{r36c:eq:neutralidad}
\end{equation}
Remainder, carry and occupation act jointly: modular reduction preserves the remainder, and the integer register distinguishes the distributions producing it.

\begin{lemma}[Resolution of regional compatibility]
\label{r36c:lem:compatibilidad} The equation $q(s)=\widetilde q(s)$ has, in $\mathbb F_3^6$, the three solutions
\[
020010,\qquad121200,\qquad222120.
\]
Its intersection with the visible repertoire is $\{121200,222120\}$.
\end{lemma}
\begin{proof}
The equation is written $sD=z$, where
\[
D=(A^{\rm o})^{-1}L_0^3(A^{\rm o})^3
   -L_0^2L_1^2A^{\rm o}L_1^3
 =\begin{pmatrix}
2&1&1&2&2&2\\0&2&0&0&0&1\\0&0&0&0&1&0\\
1&2&2&2&2&0\\1&2&2&0&1&2\\1&2&0&2&2&2
\end{pmatrix}.
\]
Elimination over $\mathbb F_3$ gives the parametrization
\[
s=(t,2,t,2t,1+2t,0),\qquad t\in\mathbb F_3.
\]
Substitution verifies the six equations and exhausts their solutions. In the generated repertoire $\mathcal W$, the two words corresponding to $t=1,2$ are present, and the one corresponding to $t=0$ is absent. This final check is an exact query of the visible-word census, preceding any Archimedean reading.
\end{proof}

\begin{proposition}[Determination of the self-scaling axis]
Compatibility, unit carry and column occupation select $s_{\rm au}=121200$ from the two visible solutions of Lemma~\ref{r36c:lem:compatibilidad}. We obtain
\[
b_4(s_{\rm au})=010200,\qquad
m(s_{\rm au})=210112,\qquad q=012120,
\]
with occupation $223232$.
\end{proposition}
\begin{proof}
For $s=222120$, we compute $b_4(s)=222000$, $m(s)=201010$ and $q(s)=000021$. In its third column, $m_3=1$ and $z_3=2$. The required occupation is three, so $x_3,y_3$ must both be nonzero. The only pair with $x_3+y_3\equiv1\pmod3$ is $(2,2)$; its integer sum with $z_3=2$ is six and produces carry two. This contradicts $c_3=1$.

For $s=121200$, the six columns respectively admit the pairs
\[
\begin{array}{c|cccccc}
j&1&2&3&4&5&6\\\hline
(x_j,y_j)&(0,2),(2,0)&(2,2)&(1,2),(2,1)&(2,2)&(2,2)&(0,2),(2,0).
\end{array}
\]
Each pair satisfies the remainder, carry and occupation equations. There are therefore eight arrangements before dual neutrality is applied.
\end{proof}

\begin{theorem}[Exact enumeration of boundary orientations]
\label{r36c:thm:dos} The preceding conditions determine exactly
\[
B_- =\begin{pmatrix}021222\\222220\\102011\end{pmatrix},
\qquad
B_+ =\begin{pmatrix}222220\\021222\\102011\end{pmatrix}.
\]
Both have remainder $012120$, carry $111111$ and occupation $223232$.
\end{theorem}
\begin{proof}
The preceding proof reduces enumeration to eight possibilities. Since $A^{\rm o}\mathbf1=(2,1,1,1,1,1)^{\mathsf T}$, neutrality of a row $v$ is equivalent to $2v_1+v_2+\cdots+v_6=0$ in $\mathbb F_3$. Its evaluation on the eight arrangements is
\[
\begin{array}{c|c|c}
x&y&(xA^{\rm o}\mathbf1,yA^{\rm o}\mathbf1,zA^{\rm o}\mathbf1)\\\hline
021220&222222&(1,2,0)\\
021222&222220&(0,0,0)\\
022220&221222&(2,1,0)\\
022222&221220&(1,2,0)\\
221220&022222&(2,1,0)\\
221222&022220&(1,2,0)\\
222220&021222&(0,0,0)\\
222222&021220&(2,1,0).
\end{array}
\]
The second and seventh rows are the only solutions. The integer column sums of either are $(3,4,5,4,5,3)$, recovering the stated remainder and carry.
\end{proof}

\subsection{Directed length and boundary orientation}

Regional transport defines a directed graph on $\mathbb F_3^6$, with edges $v\to vL_0$ and $v\to vL_1$. For mutually reachable vectors $u,v$, let
\[
d(u,v)=\min\{n\ge0:\exists\epsilon_1,\ldots,\epsilon_n\in\{0,1\},
\ v=uL_{\epsilon_1}\cdots L_{\epsilon_n}\}.
\]
This length is a combinatorial datum of transport. Its sum over the three rows preserves the regional labels:
\[
\mathcal L(B)=d(110221,B_1)+d(102011,B_2)+d(010200,B_3).
\]

\begin{proposition}[Selection by directed length]
\label{r36c:prop:minimo} We have $\mathcal L(B_-)=23$ and $\mathcal L(B_+)=20$. On the domain of the two admissible boundaries, the minimum is unique and selects $B_+$.
\end{proposition}
\begin{proof}
The shortest paths and their lengths are
\[
\begin{array}{c|c|c|c|c}
 &\text{origin}&\text{target}&\text{transport word}&d\\\hline
B_-&110221&021222&01100011&8\\
 &102011&222220&10111111&8\\
 &010200&102011&1010101&7\\\hline
B_+&110221&222220&1111&4\\
 &102011&021222&001100101&9\\
 &010200&102011&1010101&7.
\end{array}
\]
Each word is verified by successive multiplication of its matrices. To verify minimality as well, form the disjoint layers
\[
F_0(u)=\{u\},\qquad
F_{n+1}(u)=\bigl(F_n(u)L_0\cup F_n(u)L_1\bigr)
\setminus\bigcup_{j=0}^{n}F_j(u).
\]
Induction on $n$ proves that $F_n(u)$ is exactly the set of vertices at distance $n$: every shortest word of length $n+1$ comes from a shortest word of length $n$, and quotienting by previous layers eliminates shorter paths. The enumeration obtained with the two matrices displayed above gives
\[
\begin{array}{c|rrrrrrrrrr}
u\backslash n&0&1&2&3&4&5&6&7&8&9\\\hline
110221&1&2&4&8&16&32&60&105&156&177\\
102011&1&2&3&5&10&20&37&71&119&171\\
010200&1&2&4&8&15&30&56&100&148&172.
\end{array}
\]
The destinations in the first table occur for the first time in the layers indicated there. This enumeration is finite, uses only products modulo three and agrees with the breadth-first traversal of the executable certificate. Summing the three lengths gives $8+8+7=23$ and $4+9+7=20$.
\end{proof}

Orientation of the APP representation through half-open cylinders previously selects the same matrix $B_+$ at this boundary. Theorem~\ref{r36c:thm:dos} and Proposition~\ref{r36c:prop:minimo} thus establish agreement of both selections on the concrete set $\{B_-,B_+\}$. The domain of this agreement is the first joint boundary. Coinductive extension preserves its own admissibility, memory, survival and orientation relations at each depth.

\subsection{Naturality of the exceptional change of representation}

The matrix used by the exceptional reader in the active representation is
\[
A^{\rm a}=\begin{pmatrix}
0&1&1&1&1&1\\1&0&1&2&2&1\\1&1&0&1&2&2\\
1&2&1&0&1&2\\1&2&2&1&0&1\\1&1&2&2&1&0
\end{pmatrix}.
\]
The coordinate permutation
\[
R_p(v_1,v_2,v_3,v_4,v_5,v_6)
=(v_1,v_2,v_3,v_5,v_6,v_4)
\]
simultaneously transports the word and the incidence matrix.

% REMATE_LOCAL_R36
\Needspace{36\baselineskip}
\begin{theorem}[Intertwining of boundary incidence]
\label{r36c:thm:cambio}
For every $v\in\mathbb F_3^6$,
\[
R_p(vA^{\rm o})=R_p(v)A^{\rm a}.
\]
The transformation is invertible and preserves dual neutrality, remainders, carries and occupations expressed in the transported coordinates. In particular, the selected triple has the two representations
\[
\begin{array}{c|c|c}
&\text{source representation}&\text{exceptional representation}\\\hline
\text{closure}&222220&222202\\
\text{propagation}&021222&021222\\
\text{self-scaling}&102011&102110.
\end{array}
\]
\end{theorem}
\begin{proof}
If $p=(1,2,3,5,6,4)$, comparing the entries of both matrices gives $A^{\rm a}_{ij}=A^{\rm o}_{p(i),p(j)}$. For coordinate $j$,
\[
\bigl(R_p(v)A^{\rm a}\bigr)_j
=\sum_i v_{p(i)}A^{\rm o}_{p(i),p(j)}
=\sum_kv_kA^{\rm o}_{k,p(j)}
=\bigl(R_p(vA^{\rm o})\bigr)_j.
\]
The inverse is $R_p^{-1}(w)=(w_1,w_2,w_3,w_6,w_4,w_5)$. Since $R_p$ permutes coordinates, it preserves their sum and transports coordinatewise the integer sums, their remainders modulo three, quotients and numbers of nonzero entries. Preservation of all the stated registers follows. The table is obtained by applying $R_p$ to the three words of $B_+$, and the inverse recovers their original words exactly.
\end{proof}

The positional reader and incidence reader remain associated with the same enriched state through this explicit change of representation. If $\ell_{\rm pos}$ is the reader in the original representation, its expression in the transported representation is $\ell_{\rm pos}\circ R_p^{-1}$. This composition preserves positional order: changing the incidence representation and changing the reader are coordinated operations. The first boundary thus enters extension with its three words, regional orientation, integer carry and exceptional frame recoverable.
